%------------------------------------------------------------------------------%
% \chapter{Matrizes}
%------------------------------------------------------------------------------%

%------------------------------------------------------------------------------%
\begin{exercicios}{Revisão}

    \begin{exercicio}
        Resolva a equação  matricial $AX = B$, dadas as matrizes
        \[A = \left[\begin{array}{rr}
        2 & -1 \\
        -3 &  2
        \end{array} \right]	
        \qquad
        B = \left[\begin{array}{rr}
        3 &  2 \\
        -6 & -5
        \end{array} \right]	\]
    \end{exercicio}
    
\end{exercicios}

%------------------------------------------------------------------------------%
\begin{exercicios}{Desafio}
    
    \begin{exercicio}
        Determinar a inversa de cada uma das matrizes
        \begin{mathlist}[2]
        \item A = \left[\begin{array}{rr}
            5 & 6 \\ 
            4 & 5
        \end{array} \right]
        \item B = \left[\begin{array}{rr}
            2 & 5 \\ 
            1 & 3
        \end{array} \right]
        \item C = \left[\begin{array}{rr}
            1 & 0 \\ 
            0 & 2
        \end{array} \right]
        \item D = \left[\begin{array}{rr}
            1 & -1  \\ 
            1 &  1
        \end{array} \right]
        \end{mathlist}
    \end{exercicio}
    
    \begin{exercicio}
        Calcule a inversa da matriz
        \( A = \left[\begin{array}{cc}
        2 & 5 \\
        1 & 3
        \end{array} \right] \)
    \end{exercicio}
    
    \begin{exercicio}
        Um aluno, ao inverter a matriz
        \[A = \left[\begin{array}{ccc}
        1 & a & b \\
        0 & c & d \\
        4 & e & f
        \end{array} \right]
        =\left[a_{ij}\right], \qquad 1 \leq	i, j \leq 3
        \]
        cometeu um engano, e considerou o elemento $a_{31}$ igual a $3$, de forma que acabou invertendo a matriz
        \[B = \left[\begin{array}{ccc}
        1 & a & b \\
        0 & c & d \\
        3 & e & f
        \end{array} \right]	= \left[b_{ij} \right] 
        \]
        Com esse engano, o aluno encontrou 
        \[B^{-1} = \left[\begin{array}{rrr}
        \frac{5}{2} & 0 & -\frac{1}{2} \\[2mm]
        3 & 1 &           -1 \\[2mm]
        -\frac{3}{2} & 0 &  \frac{1}{2}
        \end{array} \right]\]
        Determinar $A^{-1}$.
    \end{exercicio}
    
    \begin{exercicio}
        Considere o sistema
        \[\begin{cases}
        x - my       &= 1 - m \\[2mm]
        (1 + m)x + y &= 1
        \end{cases}\]
        \begin{alphalist}
            \item Prove que o sistema admite solução única para cada número real $m$.
            \item Determine $m$ para que o valor de $x$ seja o maior possível.
        \end{alphalist}
    \end{exercicio}
    
    \begin{exercicio}
        Dada a matriz 
        \[A = \left[\begin{array}{cc}
        -2 & 3 \\
        -1 & 2
        \end{array} \right]\]
        calcule a sua inversa $A^{-1}$. 
        A relação especial, que você deve ter observado ente $A$ e $A^{-1}$ acima, 
        seria também encontrada se calculássemos as matrizes inversas de
        \begin{mathlist}
            \item \left[\begin{array}{rr}
                -3 & 4 \\ 
                -2 & 3
            \end{array} \right]
            \item \left[\begin{array}{rr}
                -5 & 6 \\ 
                -4 & 5
            \end{array} \right]
            \item \left[\begin{array}{rr}
                -1 & 2 \\
                0 & 1
            \end{array} \right]
        \end{mathlist}
        Generalize e demonstre o resultado observado.
    \end{exercicio}
    
    \begin{exercicio}
        Calcule a matriz inversa da matriz
        \[M = \left[\begin{array}{rrr}
        1 & -1 & 0 \\
        2 &  3 & 0 \\
        0 &  1 & 4
        \end{array} \right]\]
    \end{exercicio}
    
\end{exercicios}

%------------------------------------------------------------------------------%
